\documentclass[a4paper,11pt]{article}
\usepackage[pdfbookmark]{ou-tma}
\usepackage{ou-genai}

\myname{A. Student}
\mypin{X1234567}
\mycourse{M248}
\mytma{01}
\setdate{October 2026}

%% ---------------------------------------------------------------
%% Future-proofing demo: a hypothetical later presentation whose
%% proforma gained one box and lost another. The m248-26j block in
%% ou-genai.sty is untouched; this variant is built on top of it.
%% A real one would go in ou-genai.sty, but a preamble works too.
%% ---------------------------------------------------------------
\genaidefinevariant{demo-27j}{
  inherit      = m248-26j,
  presentation = {27J},
  items = {
    coding = {To write or check code}
             {e.g. an R script to draw a plot.},
  },
  remove    = {improving},
  movetoend = {none},        % new items append, so keep "did not use" last
}

\begin{document}

%% 1. The M248 26J proforma, zero configuration: module comes from
%%    \mycourse, presentation from the variant. Two boxes ticked, one
%%    carrying a note.
\genaiproforma[learning={worked through the Unit 2 derivations}, improving]

\begin{question}
  \qpart
  The proforma above is what \verb|\genaiproforma[...]| produces with the
  default variant. Everything below is the same command with other keys.
\end{question}

\clearpage

%% 2. The "did not use" case, framed, instruction line hidden.
\genaiproforma[none, frame, showinstruction=false]

\bigskip

%% 3. The demo variant: inherits m248-26j, adds `coding', drops `improving'.
%%    Module override shows the key beating \mycourse.
\genaiproforma[variant=demo-27j, module=M249, coding, starting]

\end{document}
